\section{Continuous Random Variables}

We return to the axiomatic approach. Recall that in \autoref{def:random-variable}, we required that
\[
    \brce{X \leq x} = \set{\omega}{X\para{\omega} \leq x} \in \eventSpace.
\]

Recall we defined in \autoref{def:probability-distribution-function} the Probability/Cumulative Distribution Function. These look like \autoref{fig:cdf-discrete} for discrete random variables.

\begin{figure}[ht!]
    \centering
    \begin{tikzpicture}
        \begin{axis}[
            axis lines = middle,
            xlabel = {\(x\)},
            ylabel = {\(F\para{x}\)},
            xmin = -0.25, xmax = 4.25,
            ymin = -0.25, ymax = 1.25,
            xtick = \empty,
            ytick = \empty,
            domain = 0:1.05,
            samples = 100,
            clip = false
        ]
            \addplot[no marks, thick] coordinates {(-0.2, 0) (1.2, 0)};
            \addplot[no marks, thick] coordinates {(1.2, 0.2) (1.7, 0.2)};
            \addplot[no marks, thick] coordinates {(1.7, 0.4) (2.3, 0.4)};
            \addplot[no marks, thick] coordinates {(2.3, 0.5) (2.8, 0.5)};
            \addplot[no marks, thick] coordinates {(2.8, 0.8) (3.9, 0.8)};
            \addplot[no marks, thick] coordinates {(3.9, 1) (4.2, 1)};

            \addplot[dotted] coordinates {(1.2, 0) (1.2, 0.2)};
            \addplot[dotted] coordinates {(1.7, 0.2) (1.7, 0.4)};
            \addplot[dotted] coordinates {(2.3, 0.4) (2.3, 0.5)};
            \addplot[dotted] coordinates {(2.8, 0.5) (2.8, 0.8)};
            \addplot[dotted] coordinates {(3.9, 0.8) (3.9, 1)};

            \addplot[only marks, mark=*, mark size=1.5pt] coordinates {(1.2, 0.2) (1.7, 0.4) (2.3, 0.5) (2.8, 0.8) (3.9, 1)};

            \addplot[only marks, mark=o, mark size=1.5pt] coordinates {(1.2, 0) (1.7, 0.2) (2.3, 0.4) (2.8, 0.5) (3.9, 0.8)};
        \end{axis}
    \end{tikzpicture}

    \caption{Probability Distribution Function for discrete Random Variables}
    \label{fig:cdf-discrete}
\end{figure}

Although this function is discontinuous, we note that the left limit always exists (since the function is increasing), and the function is always right-continuous.

Such functions are \emph{RCLL} (Right-Continuous with Left Limit), also called \emph{c\`{a}d\`{a}lag}. All probability distribution functions satisfy this property -- this is guaranteed from the definition.

The probability distribution function satisfies the following properties.

\begin{proposition}[Properties of Probability Distribution Function]
    \label{prop:property-probability-distribution-fn}%
    \begin{enumerate}
        \item If \(x \leq y\), then \(F\para{x} \leq F\para{y}\).
        \item For all \(a, b \in \RR\) with \(a < b\), we have
        \[
            \Prob\para{a < X \leq b} = F\para{b} - F\para{a}.
        \]
    \end{enumerate}
\end{proposition}